\paragraph{Supplement to A43: the order-three mechanism is special.}
The automorphism used in A43 cannot be supplied on every ordinary
real-multiplication target.  More precisely, let $C$ be an ordinary
genus-two curve over an algebraically closed field of characteristic
two, put $X=C^{(1)}$, and let $B$ be the locally exact differential
line bundle for $F:C\to X$.  If an automorphism $\rho$ of $C$ of
order three induces $u$ on $J_C$ with
\[
                         u^2+u+1=0,
\]
then the actual Raynaud theta support of $C$ contains a point of
exact order three.  Consequently there are ordinary finite-field
genus-two curves with the specified integral Rosati-$\mathcal O_5$
action which admit no such geometric automorphism, even after
extending constants.

\begin{proof}
Write $\rho_1=\rho^{(1)}$ and $u_1=u^{(1)}$ for the induced
automorphisms of $X$ and $J_X$.  The quotient
$X\to X/\langle\rho_1\rangle$ is a separable degree-three map
of smooth proper curves.  Its ramification is tame because the
characteristic is two.  If $\rho_1$ had no geometric fixed point,
this quotient would be etale: a nontrivial stabilizer in its
prime-order group would be the whole group.  Riemann--Hurwitz
would then give
\[
                    2=3\bigl(2g(X/\langle\rho_1\rangle)-2\bigr),
\]
which is impossible for an integer genus.  Here the exact input is
the smooth-proper separable formula and its tame different term,
Stacks Project, Lemmas~53.12.2 and~53.12.4, Tags
\texttt{0C1D} and \texttt{0C1F}; in the etale case the different
is zero \cite{StacksRH}.  Thus a fixed point $R\in X$ exists.

The required invariance is that of the actual Raynaud line bundle.
Its defining sequence is
\[
 0\longrightarrow\mathcal O_X\longrightarrow F_*\mathcal O_C
                     \longrightarrow B\longrightarrow0
\]
\cite[Section~8.1, p.~472]{OzmanPries2019}.
The natural identity $F\rho=\rho_1F$ identifies this sequence
with its pullback by $\rho_1$: both horizontal maps in the
corresponding square are isomorphisms, and their pullbacks
identify the displayed inclusion of structure sheaves.
Taking the quotient therefore gives $\rho_1^*B\simeq B$.
The induced pushforward action also fixes the class of $B$.
Hence
\[
 z=[\mathcal O_X(R)\otimes B^{-1}]
     \in\Theta_C,\qquad u_1z=z,
     \qquad 3z=(u_1^2+u_1+1)z=0.
\]
Ordinarity gives $H^0(X,B)=0$, equivalently
$0\notin\Theta_C$; see the first paragraph of
\cite[Section~8.2, p.~472]{OzmanPries2019}.
Thus $z\ne0$, and its order is exactly three.

Now apply Finding~A37(2) with $S=\{3\}$.
It supplies an ordinary finite-field genus-two curve with its
specified integral Rosati-$\mathcal O_5$ action and with no
nonzero geometric $3$-primary point on its Raynaud theta support.
The conclusion just proved excludes the stated automorphism
over the algebraic closure, so no constant extension restores it.
In particular the polarized endomorphism of A43 cannot exist on
this curve: the genus-two Torelli theorem realizes a
polarization-preserving Jacobian automorphism by a curve
automorphism, absorbing the possible sign by the hyperelliptic
involution, and its faithful action on the Jacobian preserves
the order \cite[Theorem~12.1]{MilneJV}.
\end{proof}

The extra relation used in A43 is
$\pi+V=2T+1$ for the specific $\mathbf F_2$ model, in addition
to $\pi V=2$ and $T^2+T=1$.  For a general finite-field
target whose specified real multiplication is defined over the
ground field, $T$ commutes with the $q$-Frobenius $\pi$ and
with $V=\pi^\dagger$.  Real multiplication alone then gives only
\[
 (T-\pi)^\dagger(T-\pi)=T^2-T(\pi+V)+q,
 \qquad
 (T-\pi)^2+(T-\pi)+1
                =\pi^2-(2T+1)\pi+2.
\]
Neither expression has the special value required in A43 just
from the real-multiplication relation.  The preceding argument
shows that the missing automorphism is actually absent on some
targets, rather than merely absent from the proof.
This limits that particular inverse-family calculation; it
does not disprove another rank-one realization, a different
choice of cyclic cover, or point-section separation.
